\subsection{逆伴随表示}

设 U 为 G 在 E 中的分别连续线性表示。令 \(E'\) 为 E 的对偶。映射 \(s \mapsto {}^{t}U(s)\) 是 \(E'\) 中群 \(G^{0}\) 的线性表示；称此表示为 U 的转置。映射 \(s \mapsto {}^{t}U(s^{-1}) = {}^{t}U(s)^{-1}\) 是 G 在 \(E'\) 中的线性表示，称为 U 的逆伴随表示。

引理 3. --- 设 X 为局部紧空间，Y 和 Z 为拓扑空间，\(\varphi\) 为从 \(X \times Y\) 到 \(Z\) 的连续映射，\(\varphi_x\) 为映射 \(y \mapsto \varphi(x, y)\)，从 Y 到 \(Z\)。在空间 \(\mathcal{C}(Y), \mathcal{C}(Z)\) 上赋予紧收敛拓扑时，映射 \((x, f) \mapsto f \circ \varphi_x\) 从 \(X \times \mathcal{C}(Z)\) 到 \(\mathcal{C}(Y)\) 连续。

显然只需考虑 \(\mathbf{X}\) 为紧空间的情形。取 \((x_0, f_0) \in \mathbf{X} \times \mathcal{C}(\mathbf{Z})\)、\(\mathbf{K}\) 为 \(\mathbf{Y}\) 的紧子集，并取 \(\varepsilon > 0\)。令 \(\mathbf{K}' = \varphi(\mathbf{X} \times \mathbf{K})\)。由于 \(f_0 \circ \varphi\) 在 \(X \times K\) 上一致连续，存在邻域 \(W\)，它是 \(x_0\) 的邻域，使得 \(|f_0(\varphi(x, y)) - f_0(\varphi(x_0, y))| \leqslant \varepsilon\) 对 \(x \in W\)、\(y \in K\) 成立。另一方面，若取 \(f \in \mathcal{C}(Z)\) 使得 \(|f(z) - f_0(z)| \leqslant \varepsilon\) 对所有 \(z \in K'\) 成立，则有 \(|f(\varphi(x, y)) - f_0(\varphi(x, y))| \leqslant \varepsilon\) 对 \(x \in X\)、\(y \in K\) 成立，从而有 \(|f(\varphi(x, y)) - f_0(\varphi(x_0, y))| \leqslant 2\varepsilon\) 当 \(x \in W\)、\(y \in K\) 时。引理得证。

现在回到前面的记号。

命题 3. --- (i) 若 \(U\) 分别连续，则 \({}^t U\) 在 \(\mathbf{E}'\) 上赋予弱拓扑 \(\sigma(\mathbf{E}',\mathbf{E})\) 时分别连续。

\begin{enumerate}
\def\labelenumi{(\roman{enumi})}
\setcounter{enumi}{1}
\tightlist
\item
  若 \(\mathbf{G}\) 局部紧且 \(U\) 连续，则 \({}^t U\) 在赋予紧收敛拓扑的 \(\mathbf{E}'\) 中连续。
\end{enumerate}

断言 (i) 是直接的。断言 (ii) 由引理 3 得到，其中取 \(\mathbf{X} = \mathbf{G}\)、\(\mathbf{Y} = \mathbf{Z} = \mathbf{E}\)、\(\varphi(s, x) = U(s)x\)。

